\par\addvspace{2.4ex}\noindent
\begin{minipage}{\linewidth}
\qhead{1}{2.1}{对称性}{ch02:sec:2.1}

% ---- item_id=8d12e6467fe2 seq=2.1-001 ----
\begin{q}{2.1-001}{1988数一; 880第六章基础选择6}{ch02:q:8d12e6467fe2}
(3) 设有空间区域 $\Omega_1: x^2+y^2+z^2 \le R^2, z \ge 0$ 及 $\Omega_2: x^2+y^2+z^2 \le R^2, x \ge 0, y \ge 0, z \ge 0$，则（ ）
\begin{opts}{0.48}
\opt{A}{$\displaystyle \iiint\limits_{\Omega_1} x\,\mathrm{d}v=4\iiint\limits_{\Omega_2} x\,\mathrm{d}v.$}\hfill\opt{B}{$\displaystyle \iiint\limits_{\Omega_1} y\,\mathrm{d}v=4\iiint\limits_{\Omega_2} y\,\mathrm{d}v.$}\\[0.3ex]
\opt{C}{$\displaystyle \iiint\limits_{\Omega_1} z\,\mathrm{d}v=4\iiint\limits_{\Omega_2} z\,\mathrm{d}v.$}\hfill\opt{D}{$\displaystyle \iiint\limits_{\Omega_1} xyz\,\mathrm{d}v=4\iiint\limits_{\Omega_2} xyz\,\mathrm{d}v.$}
\end{opts}
\end{q}
\end{minipage}
\markboth{2.1\quad 对称性}{2.1\quad 对称性}


\par\addvspace{2.4ex}\noindent
\begin{minipage}{\linewidth}
\qhead{1}{2.2}{直角坐标}{ch02:sec:2.2}

% ---- item_id=83094a7266c0 seq=2.2-001 ----
\begin{q}{2.2-001}{1991数一}{ch02:q:83094a7266c0}
（3）求 $\displaystyle \iiint_{\Omega}(x^2+y^2+z)\mathrm{d}v$，其中 $\Omega$ 是由曲线 $\begin{cases} y^2=2z,\\ x=0 \end{cases}$ 绕 $z$ 轴旋转一周而成的曲面与平面 $z=4$ 所围成的立体.
\end{q}
\end{minipage}
\markboth{2.2\quad 直角坐标}{2.2\quad 直角坐标}


% ---- item_id=4ef31a58ecaa seq=2.2-002 ----
\begin{q}{2.2-002}{1997数一}{ch02:q:4ef31a58ecaa}
(1) 计算 $\displaystyle I=\iiint_{\Omega}(x^2+y^2)\,\mathrm{d}v$，其中 $\Omega$ 为平面曲线 $\begin{cases} y^2=2z,\\ x=0 \end{cases}$ 绕 $z$ 轴旋转一周形成的曲面与平面 $z=8$ 所围成的区域.
\end{q}

% ---- item_id=4307044c3310 seq=2.2-003 ----
\begin{q}{2.2-003}{660数一152}{ch02:q:4307044c3310}
152 设 $\Omega$ 是由曲线
\[
\begin{cases}
y^2=2z\\
x=0
\end{cases}
\]
绕 $z$ 轴旋转一周而成的曲面与平面 $z=2$ 和 $z=8$ 所围立体，则
$\displaystyle \iiint_{\Omega}(x^2+y^2)\,\mathrm{d}v=\underline{\hspace{4em}}.$
\end{q}

% ---- item_id=c9bfdd53c2c7 seq=2.2-004 ----
\begin{q}{2.2-004}{660数一142}{ch02:q:c9bfdd53c2c7}
计算 $\displaystyle I=\iiint_{\Omega}(x^2+y^2+z^2)\,\mathrm{d}v=\underline{\hspace{4em}}$，其中 $\Omega$ 是由 $z=x^2+y^2$，$x^2+y^2=1$ 及三个坐标面围成在第一卦限内的闭区域.
\end{q}

% ---- item_id=83128050cd52 seq=2.2-005 ----
\begin{q}{2.2-005}{660数一345}{ch02:q:83128050cd52}
设 $\Omega$ 是由曲面 $z=x^2+y^2$，$y=x$，$y=0$，$z=1$ 在第一卦限所围成的区域，$f(x,y,z)$ 在 $\Omega$ 上连续，则
$\displaystyle \iiint_{\Omega} f(x,y,z)\,dv$
等于
\begin{opts}{0.48}
\opt{A}{$\displaystyle \int_0^1 dy \int_y^{\sqrt{1-y^2}} dx \int_{x^2+y^2}^{1} f(x,y,z)\,dz.$}\hfill\opt{B}{$\displaystyle \int_0^{\frac{\sqrt{2}}{2}} dx \int_y^{\sqrt{1-y^2}} dy \int_{x^2+y^2}^{1} f(x,y,z)\,dz.$}\\[0.3ex]
\opt{C}{$\displaystyle \int_0^{\frac{\sqrt{2}}{2}} dy \int_y^{\sqrt{1-y^2}} dx \int_0^1 f(x,y,z)\,dz.$}\hfill\opt{D}{$\displaystyle \int_0^{\frac{\sqrt{2}}{2}} dy \int_y^{\sqrt{1-y^2}} dx \int_{x^2+y^2}^{1} f(x,y,z)\,dz.$}
\end{opts}
\end{q}

% ---- item_id=3b44e8e07c1e seq=2.2-006 ----
\begin{q}{2.2-006}{1000第18章基础1}{ch02:q:3b44e8e07c1e}
设平面曲线 $z^2=2x$ 绕 $x$ 轴旋转一周所得空间曲面 $\Sigma$ 与平面 $x=1,x=2$ 围成的空间区域为 $\Omega$，则
$\displaystyle I=\iiint_{\Omega}\frac{1}{x^2+y^2+z^2}\,\mathrm{d}v=\underline{\hspace{4em}}.$
\end{q}

% ---- item_id=acb82b9826e2 seq=2.2-007 ----
\begin{q}{2.2-007}{880第六章综合填空4}{ch02:q:acb82b9826e2}
设 $V$ 是由 $z=\sqrt{x^2+y^2}$，$z=1$，$z=2$ 所围成的立体，计算 $\displaystyle I=\iiint_V \frac{e^z}{\sqrt{x^2+y^2}}\,\mathrm{d}x\,\mathrm{d}y\,\mathrm{d}z=\underline{\hspace{4em}}$．
\end{q}

% ---- item_id=638593034607 seq=2.2-008 ----
\begin{q}{2.2-008}{660数一140}{ch02:q:638593034607}
设 $\Omega$ 由 $x^2+\frac{y^2}{2^2}+\frac{z^2}{3^2}\leqslant 1,\ 0\leqslant z\leqslant 1$ 所确定，则 $\displaystyle \iiint_{\Omega} z^2 \mathrm{d}v=\underline{\hspace{4em}}$。
\end{q}

% ---- item_id=e2004b44a6de seq=2.2-009 ----
\begin{q}{2.2-009}{880第六章综合填空6}{ch02:q:e2004b44a6de}
设 $V$ 是由曲面 $x^{2}+y^{2}-2z^{2}=1$、平面 $z=1$ 及 $z=2$ 所围成的区域，则 $\displaystyle I=\iiint_{V} z\,dV=\underline{\hspace{4em}}$。
\end{q}

% ---- item_id=e4c1f060f4eb seq=2.2-010 ----
\begin{q}{2.2-010}{880第六章基础解答17}{ch02:q:e4c1f060f4eb}
设 $V$ 是由 $x^2+y^2+z^2=R^2$ 与 $x^2+y^2+(z-R)^2=R^2$ 所围的区域，计算 $\displaystyle I=\iiint_V z^2\,dV$.
\end{q}

\par\addvspace{2.4ex}\noindent
\begin{minipage}{\linewidth}
\qhead{1}{2.3}{球坐标}{ch02:sec:2.3}

% ---- item_id=a2925ae7e4be seq=2.3-001 ----
\begin{q}{2.3-001}{1989数一; 660数一141}{ch02:q:a2925ae7e4be}
计算三重积分 $\displaystyle \iiint_{\Omega}(x+z)\,\mathrm{d}v$，其中 $\Omega$ 是由曲面 $z=\sqrt{x^2+y^2}$ 与 $z=\sqrt{1-x^2-y^2}$ 所围成的区域.
\end{q}
\end{minipage}
\markboth{2.3\quad 球坐标}{2.3\quad 球坐标}


% ---- item_id=a138478ad942 seq=2.3-002 ----
\begin{q}{2.3-002}{2009数一}{ch02:q:a138478ad942}
设 $\Omega=\{(x,y,z)\mid x^{2}+y^{2}+z^{2}\leqslant 1\}$，则 $\displaystyle\iiint_{\Omega} z^{2}\,\mathrm{d}x\mathrm{d}y\mathrm{d}z=\underline{\hspace{4em}}$
\end{q}

% ---- item_id=1b89cecfe9b7 seq=2.3-003 ----
\begin{q}{2.3-003}{2026数一}{ch02:q:1b89cecfe9b7}
已知有界区域 $\Omega$ 由曲面 $z=\sqrt{4-x^2-y^2}$ 与 $z=\sqrt{x^2+y^2}$ 围成，函数 $f(u)$ 连续，则
$\displaystyle \iiint_{\Omega} f(x^2+y^2+z^2)\,\mathrm{d}x\,\mathrm{d}y\,\mathrm{d}z = \quad (\quad)$
\begin{opts}{0.97}
\opt{A}{$\displaystyle \int_0^{2\pi} \mathrm{d}\theta \int_0^2 \mathrm{d}r \int_r^{\sqrt{4-r^2}} f(r^2+z^2) r \,\mathrm{d}z$}\\[0.3ex]
\opt{B}{$\displaystyle \int_0^{2\pi} \mathrm{d}\theta \int_0^{\sqrt{2}} \mathrm{d}r \int_0^{\sqrt{4-r^2}} f(r^2+z^2) r \,\mathrm{d}z$}\\[0.3ex]
\opt{C}{$\displaystyle \int_0^{2\pi} \mathrm{d}\theta \int_0^{\frac{\pi}{4}} \mathrm{d}\varphi \int_0^2 f(r^2) r^2 \sin\varphi \,\mathrm{d}r$}\\[0.3ex]
\opt{D}{$\displaystyle \int_0^{2\pi} \mathrm{d}\theta \int_0^{\frac{\pi}{2}} \mathrm{d}\varphi \int_0^2 f(r^2) r^2 \sin\varphi \,\mathrm{d}r$}
\end{opts}
\end{q}

% ---- item_id=4e26fe80d89b seq=2.3-004 ----
\begin{q}{2.3-004}{1000第18章强化4}{ch02:q:4e26fe80d89b}
已知有界闭区域 $\Omega$ 由锥面 $z=\sqrt{x^2+y^2}$ 与平面 $z=1$ 围成，计算三重积分
\[
I=\iiint_{\Omega}(x+y+z)^2\,\mathrm{d}x\mathrm{d}y\mathrm{d}z.
\]
\end{q}

% ---- item_id=41d0a15ec674 seq=2.3-005 ----
\begin{q}{2.3-005}{880第六章基础解答14}{ch02:q:41d0a15ec674}
设 $V$ 由曲面 $z=\sqrt{R^2-x^2-y^2}$ 与 $z=\sqrt{x^2+y^2}$ 所围，求 $\displaystyle I=\iiint_V z\,dV$。
\end{q}

% ---- item_id=b77864095f6b seq=2.3-006 ----
\begin{q}{2.3-006}{880第六章基础解答15}{ch02:q:b77864095f6b}
设 $V$ 是由曲面 $z=\sqrt{1-x^{2}-y^{2}}$ 与 $z+1=\sqrt{x^{2}+y^{2}}$ 所围的区域，计算 $\displaystyle I=\iiint_{V} z^{2}\,dV$.
\end{q}

% ---- item_id=d2426596fbcd seq=2.3-007 ----
\begin{q}{2.3-007}{1000第18章强化1}{ch02:q:d2426596fbcd}
设 $\Omega=\{(x,y,z)\mid x^2+y^2+z^2\leqslant 1,x\geqslant 0,y\geqslant 0,z\geqslant 0\}$，则 $\displaystyle \iiint_{\Omega}(x^2+2y^2+3z^2)\,\mathrm{d}x\mathrm{d}y\mathrm{d}z=\underline{\hspace{4em}}$。
\end{q}

% ---- item_id=6ad6a31c0c25 seq=2.3-008 ----
\begin{q}{2.3-008}{880第六章综合解答22}{ch02:q:6ad6a31c0c25}
设 $V:0\le z\le \sqrt{R^2-x^2-y^2}(R>0)$，求 $\displaystyle I=\iiint_V (3x^2+5y^2+7z^2)dV$.
\end{q}

% ---- item_id=0ec0041d97d2 seq=2.3-009 ----
\begin{q}{2.3-009}{2003数一}{ch02:q:0ec0041d97d2}
设函数 $f(x)$ 连续且恒大于零，有
\[
F(t)=\frac{\iiint_{\Omega(t)} f(x^{2}+y^{2}+z^{2})\,\mathrm{d}v}{\iint_{D(t)} f(x^{2}+y^{2})\,\mathrm{d}\sigma},\qquad G(t)=\frac{\iint_{D(t)} f(x^{2}+y^{2})\,\mathrm{d}\sigma}{\int_{-t}^{t} f(x^{2})\,\mathrm{d}x}
\]
其中 $\Omega(t)=\{(x,y,z)\mid x^{2}+y^{2}+z^{2}\leqslant t^{2}\}$，$D(t)=\{(x,y)\mid x^{2}+y^{2}\leqslant t^{2}\}$．

（\rn{I}）讨论 $F(t)$ 在区间 $(0,+\infty)$ 内的单调性；（\rn{II}）证明：当 $t>0$ 时，$F(t)>\dfrac{2}{\pi}G(t)$．
\end{q}

% ---- item_id=92643ca65e03 seq=2.3-010 ----
\begin{q}{2.3-010}{880第六章综合填空7}{ch02:q:92643ca65e03}
设 $V$ 由曲面 $z=\sqrt{x^2+y^2}$ 与 $z=\sqrt{1-x^2-y^2}$ 所围，则
$\displaystyle I=\iiint_V (x+z)\,dV=\underline{\hspace{4em}}.$
\end{q}

\par\addvspace{2.4ex}\noindent
\begin{minipage}{\linewidth}
\qhead{1}{2.4}{带绝对值}{ch02:sec:2.4}

% ---- item_id=5430143f39d0 seq=2.4-001 ----
\begin{q}{2.4-001}{真题同源150}{ch02:q:5430143f39d0}
设区域 $\Omega=\{(x,y,z)\mid \sqrt{x^2+y^2}\le z\le 1\}$，计算积分 \[ I=\iiint\limits_{\Omega}\left|\sqrt{x^2+y^2+z^2}-1\right|\mathrm{d}V. \]
\end{q}
\end{minipage}
\markboth{2.4\quad 带绝对值}{2.4\quad 带绝对值}


% ---- item_id=79cd886a6f82 seq=2.4-002 ----
\begin{q}{2.4-002}{880第六章综合解答23}{ch02:q:79cd886a6f82}
设 $V$ 是由平面 $z=0,z=1$ 及圆柱面 $x^2+y^2=2$ 所围成的图形，计算 \[ I=\iiint_V \left|z-\sqrt{x^2+y^2}\right|\,dxdydz. \]
\end{q}

\par\addvspace{2.4ex}\noindent
\begin{minipage}{\linewidth}
\qhead{1}{2.5}{交换积分次序}{ch02:sec:2.5}

% ---- item_id=31e60f5f40c3 seq=2.5-001 ----
\begin{q}{2.5-001}{880第六章综合填空5}{ch02:q:31e60f5f40c3}
积分 $\displaystyle I=\int_0^1 dx \int_0^{1-x} dz \int_0^{1-x-z} (1-y)e^{-(1-y-z)^2} dy=\underline{\hspace{4em}}$。
\end{q}
\end{minipage}
\markboth{2.5\quad 交换积分次序}{2.5\quad 交换积分次序}


% ---- item_id=dd8e1f044cad seq=2.5-002 ----
\begin{q}{2.5-002}{2014同济大学}{ch02:q:dd8e1f044cad}
$\displaystyle I=\int_0^1 \mathrm{d}x\int_x^1 \mathrm{d}y\int_y^1 y\mathrm{e}^{1+z^4}\,\mathrm{d}z.$
\end{q}

% ---- item_id=758a506f6c6c seq=2.5-003 ----
\begin{q}{2.5-003}{真题同源150}{ch02:q:758a506f6c6c}
$\displaystyle I=\int_0^1\mathrm{d}x\int_x^1 y\,\mathrm{d}y\int_y^1\sqrt{1+z^4}\,\mathrm{d}z.$
\end{q}

\par\addvspace{2.4ex}\noindent
\begin{minipage}{\linewidth}
\qhead{1}{2.6}{应用}{ch02:sec:2.6}

\qhead{2}{2.6.1}{形心}{ch02:sec:2.6.1}

\qhead{3}{2.6.1.1}{求形心（质心）}{ch02:sec:2.6.1.1}

% ---- item_id=219e1dc29fda seq=2.6.1.1-001 ----
\begin{q}{2.6.1.1-001}{2010数一}{ch02:q:219e1dc29fda}
设 $\Omega=\{(x,y,z)\mid x^2+y^2\leqslant z\leqslant 1\}$，则 $\Omega$ 的形心的竖坐标 $\bar z=\underline{\hspace{4em}}$
\end{q}
\end{minipage}
\markboth{2.6.1.1\quad 求形心（质心）}{2.6.1.1\quad 求形心（质心）}


% ---- item_id=daa33b5b093e seq=2.6.1.1-002 ----
\begin{q}{2.6.1.1-002}{1000第18章基础2}{ch02:q:daa33b5b093e}
设 $\Omega$ 是由上半球面 $z=\sqrt{4-x^2-y^2}$ 与曲面 $x^2+y^2=3z$ 所围成的空间有界闭区域，则 $\Omega$ 的形心竖坐标 $\bar{z}= \underline{\hspace{4em}}$.
\end{q}

% ---- item_id=7c228568c488 seq=2.6.1.1-003 ----
\begin{q}{2.6.1.1-003}{2013数一}{ch02:q:7c228568c488}
设直线 $L$ 过 $A(1,0,0),B(0,1,1)$ 两点，将 $L$ 绕 $Z$ 轴旋转一周得到曲面 $\Sigma$，$\Sigma$ 与平面 $z=0,z=2$ 所围成的立体为 $\Omega$，

（I）求曲面 $\Sigma$ 的方程；（II）求 $\Omega$ 的形心坐标.
\end{q}

% ---- item_id=aa682f9df977 seq=2.6.1.1-004 ----
\begin{q}{2.6.1.1-004}{1000第18章强化23}{ch02:q:aa682f9df977}
设锥面 $\Sigma$ 的顶点是 $A(0,1,1)$，准线是 $\begin{cases} x^2+y^2=1,\\ z=0, \end{cases}$ 直线 $L$ 过顶点 $A$ 和准线上任一点 $M_1(x_1,y_1,0)$.

$\Omega$ 是 $\Sigma(0\le z\le1)$ 与平面 $z=0$ 所围成的锥体. 求:

(1) $L$ 和 $\Sigma$ 的方程；

(2) $\Omega$ 的形心坐标.
\end{q}

% ---- item_id=c53f2a2aa310 seq=2.6.1.1-005 ----
\begin{q}{2.6.1.1-005}{2019数一}{ch02:q:c53f2a2aa310}
设 $\Omega$ 是由锥面 $x^2+(y-z)^2=(1-z)^2(0\leq z\leq 1)$ 与平面 $z=0$ 围成的锥体,求 $\Omega$ 的形心坐标.
\end{q}

% ---- item_id=8353976af09b seq=2.6.1.1-006 ----
\begin{q}{2.6.1.1-006}{2000数一}{ch02:q:8353976af09b}
设有一半径为 $R$ 的球体，$P_0$ 是此球的表面上的一个定点，球体上任一点的密度与该点到 $P_0$ 距离的平方成正比（比例常数 $k>0$），求球体的重心位置．
\end{q}

% ---- item_id=d96ca005ad20 seq=2.6.1.1-007 ----
\begin{q}{2.6.1.1-007}{真题同源150}{ch02:q:d96ca005ad20}
【例8.3】设球体 $\Omega$ 的方程为 $x^2+y^2+z^2=2Rz$，球体上任意点的密度与该点到原点的距离的平方成正比，比例常数 $k>0$，求球体的质心．
\end{q}

% ---- item_id=17d9586acd9d seq=2.6.1.1-008 ----
\begin{q}{2.6.1.1-008}{660数一158}{ch02:q:17d9586acd9d}
158 设球体 $x^2+y^2+z^2 \leqslant z$ 上任一点处的密度等于该点到原点的距离的平方，则此球的质心的 $z$ 坐标为 $\underline{\hspace{4em}}$.
\end{q}

% ---- item_id=6742040d8759 seq=2.6.1.1-009 ----
\begin{q}{2.6.1.1-009}{880第六章综合解答24}{ch02:q:6742040d8759}
（24）某均匀物体由上、下两部分组成，上部分是半径为 $a$ 的半球体，下部分是底面半径为 $a$、高为 $3$ 的直圆锥体，且半球体的底面圆与圆锥的底面重合，问当 $a$ 为何值时，此物体的质心恰好在球心位置？
\end{q}

% ---- item_id=cfe14e8febf7 seq=2.6.1.1-010 ----
\begin{q}{2.6.1.1-010}{1000第17章强化5}{ch02:q:cfe14e8febf7}
求抛物面 $\Sigma:z=x^2+y^2$ 上的点到空间图形 $\Omega:x^2+y^2\leq z\leq 1$ 的形心的最短距离.
\end{q}

\par\addvspace{2.4ex}\noindent
\begin{minipage}{\linewidth}
\qhead{3}{2.6.1.2}{逆用形心公式}{ch02:sec:2.6.1.2}

% ---- item_id=8eb3e7f8c579 seq=2.6.1.2-001 ----
\begin{q}{2.6.1.2-001}{2015数一}{ch02:q:8eb3e7f8c579}
（12）设 $\Omega$ 是由平面 $x+y+z=1$ 与三个坐标平面所围成的空间区域，则
$\displaystyle \iiint_{\Omega}(x+2y+3z)\,\mathrm{d}x\mathrm{d}y\mathrm{d}z=\underline{\hspace{4em}}$
\end{q}
\end{minipage}
\markboth{2.6.1.2\quad 逆用形心公式}{2.6.1.2\quad 逆用形心公式}


% ---- item_id=41652d3a3c11 seq=2.6.1.2-002 ----
\begin{q}{2.6.1.2-002}{880第六章基础填空11}{ch02:q:41652d3a3c11}
设 $V:x^2+y^2+z^2\leqslant 2y-2x$，则 $\displaystyle I=\iiint_V (x+y+z)\,dV= \underline{\hspace{4em}}$．
\end{q}

\par\addvspace{2.4ex}\noindent
\begin{minipage}{\linewidth}
\qhead{2}{2.6.2}{实心体质量}{ch02:sec:2.6.2}

% ---- item_id=5a95b72eddc9 seq=2.6.2-001 ----
\begin{q}{2.6.2-001}{880第九章基础填空9}{ch02:q:5a95b72eddc9}
设立体 $\Omega$ 为 $x^2+y^2+z^2\le 2z$，其密度 $\rho(x,y,z)=x^2$，则 $\Omega$ 的质量为 $\underline{\hspace{4em}}$．
\end{q}
\end{minipage}
\markboth{2.6.2\quad 实心体质量}{2.6.2\quad 实心体质量}


\par\addvspace{2.4ex}\noindent
\begin{minipage}{\linewidth}
\qhead{2}{2.6.3}{实心体体积}{ch02:sec:2.6.3}

% ---- item_id=c62adb73a76d seq=2.6.3-001 ----
\begin{q}{2.6.3-001}{1994数一}{ch02:q:c62adb73a76d}
已知点 $A$ 与点 $B$ 的直角坐标分别为 $(1,0,0)$ 与 $(0,1,1)$，线段 $AB$ 绕 $z$ 轴旋转一周所成的旋转曲面为 $S$，求由 $S$ 及两平面 $z=0$，$z=1$ 所围成的立体体积。
\end{q}
\end{minipage}
\markboth{2.6.3\quad 实心体体积}{2.6.3\quad 实心体体积}


% ---- item_id=bbe0a1d738da seq=2.6.3-002 ----
\begin{q}{2.6.3-002}{880第六章基础填空12}{ch02:q:bbe0a1d738da}
球体 $x^2+y^2+z^2=R^2(R>0)$ 被圆柱面 $x^2+y^2=Rx$ 所截得含在圆柱面内的立体的体积为 $\underline{\hspace{4em}}$．
\end{q}

% ---- item_id=7a93f39e13dc seq=2.6.3-003 ----
\begin{q}{2.6.3-003}{880第六章基础解答16}{ch02:q:7a93f39e13dc}
求曲面 $z=\sqrt{5-x^2-y^2}$ 与 $x^2+y^2=4z$ 所围立体体积.
\end{q}

\par\addvspace{2.4ex}\noindent
\begin{minipage}{\linewidth}
\qhead{2}{2.6.4}{转动惯量}{ch02:sec:2.6.4}

% ---- item_id=e43f1cef91ac seq=2.6.4-001 ----
\begin{q}{2.6.4-001}{660数一159}{ch02:q:e43f1cef91ac}
设 $\Gamma$ 为质量均匀分布的半圆 $y=\sqrt{1-x^2}$，线密度为 $\rho$，则 $\Gamma$ 对 $x$ 轴的转动惯量 $I_x=\underline{\hspace{4em}}$．
\end{q}
\end{minipage}
\markboth{2.6.4\quad 转动惯量}{2.6.4\quad 转动惯量}


% ---- item_id=d78bb8d38587 seq=2.6.4-002 ----
\begin{q}{2.6.4-002}{880第六章基础填空14}{ch02:q:d78bb8d38587}
(14) 设平面薄片（密度 $\rho=1$）由 $y^2=x^3$ 与直线 $y=x$ 所围，则 $D$ 对 $x$ 轴和 $y$ 轴的转动惯量分别为 $I_x=\underline{\hspace{4em}}$，$I_y=\underline{\hspace{4em}}$。
\end{q}
